\section*{写给谁，怎么读}
\phantomsection\label{front:reading-routes}
\begingroup
% Count the union of actual page spans after cross-references settle. Keeping
% the ranges here lets the PDF and its navigation audit use the same selection.
\ExplSyntaxOn
\seq_new:N \l__reading_pages_seq
\NewDocumentCommand{\setroutebudget}{mm}{
  \group_begin:
  \seq_clear:N \l__reading_pages_seq
  \clist_map_inline:nn {#2}{
    \seq_set_split:Nnn \l_tmpa_seq {/}{##1}
    \int_set:Nn \l_tmpa_int {\getpagerefnumber{\seq_item:Nn \l_tmpa_seq {1}}}
    \int_set:Nn \l_tmpb_int {\getpagerefnumber{\seq_item:Nn \l_tmpa_seq {2}}}
    \int_compare:nNnT {\l_tmpa_int} > {0}{
      \int_step_inline:nnn {\l_tmpa_int}{\l_tmpb_int}{
        \seq_put_right:Nn \l__reading_pages_seq {####1}
      }
    }
  }
  \seq_remove_duplicates:N \l__reading_pages_seq
  \cs_gset:cpx {readingbudget#1}{\seq_count:N \l__reading_pages_seq}
  \group_end:
}
\NewDocumentCommand{\routebudget}{m}{\use:c {readingbudget#1}}
\NewExpandableDocumentCommand{\routevisits}{m}{
  \int_div_truncate:nn {\use:c {readingbudget#1}+4}{5}
}
\ExplSyntaxOff
\setroutebudget{quick}{intro:00-functions/route:gradient-end,intro:07-functions-relu/route:relu-end,intro:17-backpropagation/route:backprop-end,intro:01-vectors/route:mlp-forward-end,sec:mlp-ffn/route:mlp-end,intro:20a-tokenizer/route:joint-core-end,intro:23-llm-systems/route:generation-end,intro:23aa-pretraining-icl/route:icl-end,intro:23d-agent-systems/route:protocols-end}
\setroutebudget{standard}{intro:00-functions/route:modern-core-end,intro:23e-tech-judgement/route:standard-end}
\setroutebudget{hands}{intro:00-functions/route:backprop-all-end,intro:01-vectors/route:mlp-end,intro:20b-task-contract/route:embedding-end,intro:21a-context-gap/route:joint-all-end,intro:23aa-pretraining-icl/route:feedback-end,app:rational-mini/route:lab-end}
\setlength{\parskip}{5pt}
\noindent\textbf{适合谁。}本书写给想弄清“语言模型和 AI 助手到底怎样计算”的读者。会解简单方程、代入函数，就可以开始阅读。不要求先学过线性代数、概率或深度学习；对数、导数等运算生疏时，可以结合正文例子查阅数学速查。选择动手路线，还需要能运行 Python 脚本。

\noindent\textbf{读完能做到什么。}复算一次从输入到损失、反向传播和参数更新的完整循环；解释 Q/K/V、因果 Mask、FFN 与输出头的分工；从模型配置估算 KV Cache；区分模型生成、检索、工具执行与权限检查的职责。

\tagpdfsetup{table/header-rows={1}}
\begin{center}\small
\begin{tabularx}{\linewidth}{@{}p{29mm}p{44mm}X@{}}
\toprule
路线与篇幅 & 适合谁 & 完成后的主要收获\\\midrule
速读，约\routebudget{quick}页 & 先想建立整体认识 & 解释学习、上下文读取、生成与助手系统怎样连接\\
标准，约\routebudget{standard}页 & 希望理解并复算主要机制 & 算通训练循环，理解 Transformer，并从模型配置估算缓存\\
动手，约\routebudget{hands}页 & 愿意运行代码、核对数字 & 复现前向计算、梯度和参数更新，验证实验结果\\
\bottomrule
\end{tabularx}
\end{center}
\noindent\textbf{怎样安排时间。}按每次读 5 页安排，速读约需\routevisits{quick}次，标准约需\routevisits{standard}次。
这只是排期示例；页数按选定段落覆盖的整页去重统计，练习、补数学与运行脚本的时间另计。
动手路线包含手算和代码实验，耗时还取决于复算进度。
各路线的具体范围见\hyperref[app:reading-details]{附录~\ref*{app:reading-details} 的阅读细目（第~\pageref*{app:reading-details}~页）}。

对数学基础没把握，可以先做\hyperref[app:math-entry]{五题自检（第~\pageref*{app:math-entry}~页）}；
也可以直接进入序章，遇到生疏运算时再补读。

\noindent\textbf{内容范围。}主线用小模型展示完整计算，另有真实模型的配置与运行对照。训练、微调和部署侧重原理与估算，不提供生产工程的完整操作教程；视觉扩展介绍图像输入，不展开图像、音频生成，也不系统展开拒答与安全对齐训练。

选好路线，就从下一页的序章开始。
\endgroup
\phantomsection\label{front:reading-end}
